\section{Discussion}

\subsection{CPU Pipeline Limitations}

The CPU-based pipeline demonstrates strong efficiency at moderate input resolutions but becomes increasingly constrained as image resolution grows. Since feature detection and descriptor extraction scale with the number of pixels, higher resolutions introduce significant computational overhead even with SIMD vectorization. 

Reducing camera resolution improves runtime performance but introduces trade-offs in feature quality. At lower resolutions, distant objects occupy fewer pixels, leading to weaker or less distinctive descriptors. This reduces matching reliability, particularly under multiscale conditions, and can negatively impact tracking robustness.

\subsection{GPU Pipeline Trade-offs}

The GPU-accelerated pipeline alleviates many of the computational bottlenecks associated with feature extraction and descriptor computation by exploiting data-parallel execution. However, its performance remains dependent on memory transfer costs and overall GPU availability.

Although the use of storage buffers significantly reduces data transfer overhead compared to transferring full image buffers, the cost of transferring match data still scales with the number of correspondences. As a result, scene complexity directly influences runtime performance.

Additionally, GPU performance on mobile platforms is subject to system-level constraints. Since the GPU is often shared with rendering and other application processes, scheduling variability and resource contention can reduce the effective throughput available for feature processing.

\subsection{Geometric Tracking Advantages}

The proposed $\mathrm{SE}(3)$ render-and-track formulation addresses several limitations of purely algebraic pipelines by incorporating explicit geometric constraints derived from the object model. The availability of object-space coordinates during rendering enables direct lifting of image correspondences to 2D--3D constraints, reducing sensitivity to clustered or spatially limited feature distributions.

Furthermore, the tension-based refinement introduces adaptive correspondence weighting within the optimization process. By assigning weights based on relative residual magnitudes, the system emphasizes inconsistent correspondences while suppressing those already in agreement. This results in improved robustness under partial visibility, noise, and viewpoint changes, without requiring manually tuned robust loss functions.

\subsection{Comparison with ARCore}

ARCore provides a highly optimized mobile tracking framework with built-in image tracking capabilities \cite{arcore}. While it achieves strong performance in marker-based scenarios, it operates as a closed system with limited control over the underlying tracking pipeline. Consequently, its performance may degrade under strong viewpoint changes or partial occlusions, depending on the visual properties of the tracked image.

In contrast, the proposed system offers explicit control over all stages of the tracking pipeline and enables the integration of geometry-aware constraints through render-and-track optimization. This flexibility allows the system to maintain stable pose estimates in scenarios where purely appearance-based tracking methods struggle.

\subsection{Future Work}

While the proposed framework focuses on feature-based tracking using keypoint correspondences, several extensions remain possible. One promising direction is the integration of silhouette-based tracking methods. Keypoint-based approaches rely on the presence of distinctive local features, which may become sparse under low-texture conditions, motion blur, or significant scale variation.

Silhouette-based tracking, which aligns projected object contours with observed image edges, can provide complementary geometric constraints that are less dependent on texture. Incorporating silhouette alignment within the render-and-track framework could improve robustness in challenging scenarios, particularly under partial visibility or low-texture surfaces.

Another direction for future work is the exploration of hybrid optimization strategies that combine feature-based residuals with additional geometric cues. Extending the tension-based formulation to incorporate multiple constraint types may further improve robustness while maintaining computational efficiency.

Investigating these extensions within the proposed framework represents a promising direction for improving real-time tracking performance on resource-constrained devices.